\section{Introduction}

Classical D-stability asks whether the spectrum of a real matrix remains in the open left half-plane after multiplication by every positive diagonal matrix. In dimension three this problem admits a complete description in terms of principal minors and an exact homogeneous minimization, due to Cain \citep{Cain1976}; Bahl and Cain subsequently characterized fixed inertia under all positive diagonal multipliers \citep{BahlCain1977}. Generalized matrix-stability frameworks extend this question from the left half-plane to more general spectral regions \citep{Kushel2019}, and Kushel and Pavani developed an abstract forbidden-boundary characterization for positive-diagonal stability in conic regions and their complements \citep{KushelPavani2022}.

For a commensurate Caputo system of order \(0<\alpha<1\), asymptotic stability is governed by the Matignon angular region
\[
\Sigma_\alpha
=
\{z\neq0:|\arg z|>\alpha\pi/2\},
\]
rather than by the left half-plane \citep{Matignon1996}; see also the modern treatment in \citet{BrandiburGarrappaKaslik2021}. Fixed-polynomial fractional root-location criteria are well developed \citep{CermakNechvatal2017,Bourafa2020}, and structured cyclic networks admit fractional secant conditions \citep{Siami2021}. In closely related ecological work, Alraddadi and Alharthi derived exact order-dependent stability atlases for a Caputo Allee predator--prey model and isolated a region of purely fractional stabilization, where the fractional equilibrium is stable despite a right-half-plane Jacobian spectrum \citep{AlraddadiAlharthi2026}. That model-specific phenomenon motivates, but does not resolve, the positive-diagonal orbit problem studied here. What is not supplied by these results is an explicit elimination of the universal positive-diagonal multiplier for a generic real \(3\times3\) matrix in the full Matignon reflex sector.

The purpose of this paper is to solve that low-dimensional problem exactly on the full-dimensional strict-\(P(-A)\) stratum. We show that the entire positive-diagonal orbit is determined by four invariants,
\[
(\beta_{12},\beta_{13},\beta_{23},\kappa),
\]
and that the all-diagonal Matignon condition reduces to a single scalar inequality
\[
\kappa<T_\alpha(\beta).
\]
The threshold \(T_\alpha\) is defined by a two-dimensional simplex minimization of an exact cubic boundary function. At \(\alpha=1\) it collapses to Cain's classical threshold, while for \(2/3<\alpha<1\) it lies strictly above the classical surface and therefore exposes an exact genuinely fractional band.

Our main contributions are as follows.
\begin{enumerate}
\item We give an exact \(2\times2\) classification and prove that the genuinely fractional difference from classical Hurwitz D-stability has empty full-dimensional interior in dimension two.
\item We derive a necessary-and-sufficient real-\(3\times3\) criterion
\[
A\in\mathcal F_\alpha^{(3)}
\iff
\kappa<T_\alpha(\beta)
\]
on the strict-\(P(-A)\) stratum.
\item We prove that the threshold objective is strictly convex in logit simplex coordinates. Hence the optimizer is unique and nondegenerate, the threshold depends smoothly on its parameters, and the full threshold surface admits an explicit parametrization.
\item We prove that dimension three is the smallest real dimension in which the genuinely fractional positive-diagonal Matignon-stable class outside classical Hurwitz D-stability has nonempty full-dimensional interior.
\item We derive sharp asymptotics at both \(\alpha\to1^-\) and \(\alpha\downarrow2/3\).
\item For three-species generalized Lotka--Volterra systems, we identify the orbit invariants with normalized reciprocal-pair and directed three-cycle feedback coordinates.
\end{enumerate}

The contribution should be read as an explicit low-dimensional solution inside a known generalized-D-stability framework, not as the introduction of fractional D-stability itself. This distinction is essential for comparison with \citet{KushelPavani2022}, while the limit \(\alpha=1\) makes the relation to \citet{Cain1976} exact rather than merely qualitative. Relative to the preceding ecological stability-atlas study \citep{AlraddadiAlharthi2026}, the present result moves from a fixed nonlinear model and its parameter atlas to an exact matrix-orbit theorem quantified over every positive diagonal multiplier.

The phrase \emph{genuinely fractional} is used throughout in a precise orbit-level sense: a matrix is fractionally stable under every positive diagonal multiplier but is not classically Hurwitz D-stable. This is stronger than exhibiting a single matrix whose fractional system is stable while the corresponding integer-order system at \(D=I\) is unstable.

The paper is organized as follows. Section~\ref{sec:framework} introduces the generalized-D-stability setting and the projective simplex reduction of the positive diagonal orbit. Section~\ref{sec:dimension-two} settles the exact two-dimensional case. Section~\ref{sec:exact-3x3} proves the real-\(3\times3\) threshold theorem and the complete full-dimensional interior classification. Section~\ref{sec:threshold-geometry} develops the convex geometry, parametrization, asymptotics, closed-form sufficient certificate, realizability theorem, and minimum-dimension consequence. Section~\ref{sec:ecology} translates the invariant coordinates into three-species generalized Lotka--Volterra feedback loops. Section~\ref{sec:computation} reports the independent adversarial and interval-certified validation campaign, and Section~\ref{sec:discussion} places the result relative to its closest classical and generalized-D-stability predecessors.
